% Part: first-order-logic
% Chapter: natural-deduction
% Section: quantifier-rules

\documentclass[../../../include/open-logic-section]{subfiles}

\begin{document}

\olfileid{fol}{ntd}{qrl}

\olsection{Aturan Kuantor}

\subsection{Aturan kanggo $\lforall$}

\begin{defish}
\AxiomC{$!A(a)$}
\RightLabel{\Intro{\lforall}}
\UnaryInfC{$\lforall[x][\Atom{!A}{x}]$}
\DisplayProof
\hfill
\AxiomC{$\lforall[x][\Atom{!A}{x}]$}
\RightLabel{\Elim{\lforall}}
\UnaryInfC{$!A(t)$}
\DisplayProof
\end{defish}

Ing aturan kanggo~$\lforall$, $t$ iku term tertutup (term kang ora
ngemot variabel), lan $a$~iku !!a{constant} kang ora muncul ing
kesimpulan~$\lforall[x][!A(x)]$, utawa ing sembarang asumsi kang
!!{undischarged} ing !!{derivation} kang mungkasi nganggo
premis~$!A(a)$. Kita ngarani $a$ \emph{eigenvariabel} saka inferensi
\Intro{\lforall}.\footnote{Kita nggunakake istilah ``eigenvariabel''
senajan $a$ ing aturan ndhuwur iku konstanta. Iki amarga alasan
sajarah.}

\subsection{Aturan kanggo $\lexists$}

\begin{defish}
\AxiomC{$\Atom{!A}{t}$}
\RightLabel{\Intro{\lexists}}
\UnaryInfC{$\lexists[x][\Atom{!A}{x}]$}
\DisplayProof
\hfill
\AxiomC{$\lexists[x][\Atom{!A}{x}]$}
\AxiomC{[$\Atom{!A}{a}$]$^n$}
\DeduceC{$!C$}
\DischargeRule{\Elim{\lexists}}{n}
\BinaryInfC{$!C$}
\DisplayProof
\end{defish}

Maneh, $t$ iku term tertutup, lan $a$ iku !!a{constant} kang ora
muncul ing premis $\lexists[x][!A(x)]$, ing kesimpulan~$!C$, utawa ing
sembarang asumsi kang !!{undischarged} ing !!{derivation}s kang
mungkasi nganggo premis loro mau (kajaba asumsi $!A(a)$). Kita
ngarani $a$ \emph{eigenvariabel} saka inferensi \Elim{\lexists}.

Syarat yen eigenvariabel ora muncul ing kesimpulan utawa asumsi kang
!!{undischarged} ing !!{derivation}s premis kanggo inferensi
\Intro{\lforall}, lan ora muncul ing premis eksistensial, kesimpulan,
utawa asumsi kaya mangkono kanggo inferensi \Elim{\lexists}, diarani
\emph{syarat eigenvariabel}. Cathetan koreksi sumber OLPL-009:
ringkesan Inggris nyebut kabeh premis, sanajan premis universal lan
asumsi kang dibubarake mesthi ngemot eigenvariabel; rumusan ing kene
ngetutake syarat rinci kang diwenehake langsung ing ndhuwur.

\begin{explain}
Elinga konvensi yen $!A$ iku !!a{formula} kanthi !!{variable}~$x$
bebas, kita nuduhake iki kanthi nulis~$!A(x)$. Ing konteks kang padha,
$!A(t)$ banjur minangka cekakan saka~$\Subst{!A}{t}{x}$. Dadi, kita
uga bisa nulis aturan $\Intro\lexists$ minangka:
\begin{prooftree}
\AxiomC{$\Subst{!A}{t}{x}$}
\RightLabel{\Intro{\lexists}}
\UnaryInfC{$\lexists[x][!A]$}
\end{prooftree}
Gatekna yen $t$ bisa wae wis muncul ing~$!A$; tuladhane, $!A$~bisa
awujud~$\Atom{\Obj P}{t,x}$. Mula, nginferensi
$\lexists[x][\Atom{\Obj P}{t,x}]$ saka~$\Atom{\Obj P}{t,t}$ iku
panrapan $\Intro\lexists$ kang bener---kita kena ``ngganti'' siji
utawa luwih, lan ora kudu kabeh, kedadeyane~$t$ ing premis nganggo
!!{variable}~$x$ kang kaiket. Nanging, syarat eigenvariabel ing
$\Intro\lforall$ lan~$\Elim\lexists$ mbutuhake supaya !!{constant}~$a$
ora muncul ing~$!A$. Mula, kita ora bisa kanthi bener nginferensi
$\lforall[x][\Atom{\Obj P}{a,x}]$ saka $\Atom{\Obj P}{a,a}$
nganggo~$\Intro\lforall$.
\end{explain}

\begin{explain}
Ing \Intro{\lexists} lan \Elim{\lforall} ora ana watesan tambahan,
lan term tertutup~$t$ bisa dipilih sawenang-wenang, mula kita ora
perlu nguwatirake syarat liyane. Nanging, ing aturan
\Elim{\lexists} lan \Intro{\lforall}, syarat eigenvariabel mbutuhake
supaya !!{constant}~$a$ ora muncul ing panggonan kang dilarang ing
kesimpulan, premis eksistensial, utawa asumsi kang !!{undischarged},
miturut aturan sing cocog. Syarat iki perlu kanggo njamin yen sistem
iku sah, yaiku sistem mung !!{derive}s !!{sentence}s saka asumsi kang
!!{undischarged} nalika ukara kasebut pancen ndherek saka asumsi mau.
Tanpa syarat iki, inferensi iki bakal diidinake:
\begin{prooftree}
  \AxiomC{$\lexists[x][!A(x)]$}
  \AxiomC{$\Discharge{!A(a)}{1}$}
  \RightLabel{*\Intro{\lforall}}
  \UnaryInfC{$\lforall[x][!A(x)]$}
  \RightLabel{\Elim{\lexists}}
  \BinaryInfC{$\lforall[x][!A(x)]$}
\end{prooftree}
Nanging, $\lexists[x][!A(x)] \Entails/ \lforall[x][!A(x)]$.

Amarga aturan eliminasi kanggo kuantor mung ngidini substitusi term
tertutup kanggo !!{variable}s, saben !!{formula} kang bisa diasilake
saka sawijining himpunan !!{sentence}s mesthi uga !!a{sentence}.
\end{explain}


\end{document}
